\begin{textsample}{NEG Dim 2 – 2015\_c4 – Score 11.00 – 2015\_c4/enade\_2015\_secr\_exec\_875f6cdd83.txt}  \label{ex:f2_neg_350}
CAPÍTULO 4 PERCEPÇÃO DA PROVA As análises feitas neste capítulo tratam das percepções dos concluintes da Área de Secretariado Executivo sobre a prova aplicada no Enade/2015.
Essas percepções foram mensuradas por meio de nove questões que avaliaram desde o grau de dificuldade da prova até o tempo gasto para concluí -la.
As percepções sobre a prova foram relacionadas com o desempenho dos estudantes e com a Grande Região de funcionamento do curso.
O questionário de percepção da prova encontra -se ao final do Anexo VII, que traz a reprodução da prova.
O desempenho dos estudantes foi classificado em quatro quartos.
Para tanto, esse desempenho foi ordenado de forma ascendente.
O percentil 25, P25, também conhecido como primeiro quartil, é a nota de desempenho que deixa um quarto ( 25 \% ) dos valores observados abaixo e três quartos acima.
A Figura 1 apresenta uma ilustração deste conceito.
O quarto inferior de desempenho é composto pelas notas abaixo do primeiro quartil.
Já o percentil 75, P75, também conhecido como terceiro quartil, é o valor para o qual há três quartos ( 75 \% ) dos dados abaixo e um quarto acima dele.
O quarto superior de desempenho é composto pelas notas iguais ou acima do terceiro quartil.
O percentil 50, P50, também conhecido como mediana, é o valor que divide as notas em dois conjuntos de igual tamanho.
O segundo quarto inclui valores entre o primeiro quartil ( P25 ) e a mediana.
O terceiro quarto contém os valores entre a mediana ( P50 ) e o terceiro quartil ( P75 ).
Vale ressaltar que percentis, quartis e medianas são pontos que não obrigatoriamente pertencem ao conjunto original de dados, ao passo que os quartos são subconjuntos dos dados originais.
A seguir, serão apresentados gráficos com resultados selecionados, relativos às nove questões avaliadas por grupos de estudantes.
Os gráficos apresentam nas barras o percentual de alunos que assinalaram uma das opções ou a soma das porcentagens daqueles que assinalaram duas ( ou três ) delas.
Por exemplo, para as questões 1 e 2, os gráficos apresentam a porcentagem total de participantes que assinalaram as opções ( D ) Difícil e ( E ) Muito difícil.
Em cada barra foram assinalados também os extremos do intervalo de confiança de 95 \% como linhas verticais unidas por uma linha horizontal na forma da letra H maiúscula.
As Tabelas no Anexo II apresentam os valores absolutos e a distribuição percentual17 das alternativas válidas das nove questões, segundo o mesmo recorte de desempenho dos alunos e Grande Região de funcionamento do curso. 4.1 GRAU DE DIFICULDADE DA PROVA 4.1.1 Componente de Formação Geral Ao avaliarem “ Qual o grau de dificuldade desta prova na parte de Formação Geral? ” ( Questão 1 ), 20,2 \% do grupo de inscritos e presentes optaram pelas alternativas Difícil ou Muito difícil.
Entretanto, para mais da metade dos estudantes ( 68,8 \% ), o Componente de Formação Geral da prova foi considerado com grau de dificuldade Médio ( Gráfico 4.1, Gráfico 4.2 e, no Anexo II, a Tabela II.1 ).
O percentual de estudantes que consideraram a prova como Difícil ou Muito difícil foi maior na região Centro-Oeste, onde a proporção foi de 25,8 \%, enquanto a de menor incidência foi a Norte, com 15,3 \%.
No Gráfico 4.1, é possível observar que essa diferença não é estatisticamente significativa, assim como a diferença entre cada uma das regiões.
Nas Grandes Regiões, a proporção de presentes à prova que consideraram o Componente de Formação Geral como sendo de grau de dificuldade Médio esteve entre 71,3 \% na região Nordeste e 67,3 \% na região Sudeste. 17 Cumpre lembrar uma das convenções para tabelas numéricas ( pág. iii ) sobre a possibilidade de a soma das partes não resultar em 100 \% por questões de arredondamento.
O percentual de alunos que consideraram a prova Difícil ou Muito difícil foi decrescente em função dos quartos de desempenho : 24,3 \% no primeiro e no segundo quartos e 12,9 \% no quarto, grupo de melhor desempenho na prova.
No terceiro quarto de desempenho, a proporção de alunos que consideraram a prova Difícil ou Muito difícil foi de 19,6 \%.
As diferenças entre os dois primeiros quartos e o superior são estatisticamente significativas.
Para todos os quartos de desempenho, a alternativa modal para esta pergunta foi Médio, com 68,5 \% dos respondentes no primeiro quarto, 66,8 \% no segundo, 70,4 \% no terceiro e 69,3 \% no quarto superior. 4.1.2 Componente de Conhecimento Específico Ao responderem à Questão 2 – “ Qual o grau de dificuldade desta prova na parte de Componente Específico? ” – 19,8 \% do grupo de estudantes classificaram -na como Difícil ou Muito difícil.
O Componente de Conhecimento Específico da prova foi considerado com grau de dificuldade Médio por 68,5 \% dos alunos ( Gráfico 4.3, Gráfico 4.4, e, no Anexo II, a Tabela II.2 ).
A análise das respostas dos estudantes quanto ao grau de dificuldade do Componente de Conhecimento Específico da prova, agregado por Grande Região, mostra que as diferenças entre as proporções de alunos que a avaliaram como Difícil ou Muito difícil em cada uma das regiões não é estatisticamente significativa.
A maior proporção foi obtida na região Centro-Oeste ( 23,0 \% ) e a menor proporção na região Sudeste ( 17,1 \% ).
O percentual de alunos que classificaram o grau de dificuldade como Médio, no Componente de Conhecimento Específico, variou de 60,7 \% a 70,3 \%, para as regiões Norte e Nordeste, respectivamente.
Considerando -se a avaliação da dificuldade das questões do Componente de Conhecimento Específico da prova, de acordo com o desempenho dos estudantes, observa -se que não há diferença estatisticamente significativa entre as proporções de cada um dos quartos desempenho.
A proporção dos que classificaram a parte específica como Difícil ou Muito difícil variou de 17,7 \% ( último quarto ) a 22,0 \% ( segundo quarto ).
A alternativa modal para a Questão 2 foi o grau médio, com 71,6 \% do quarto inferior e 66,6 \% do superior optando por essa resposta.
O percentual do segundo quarto foi 66,8 \%, variando muito pouco em relação ao quarto superior. 4.2 EXTENSÃO DA PROVA EM RELAÇÃO AO TEMPO TOTAL Indagados quanto à extensão da prova, em relação ao tempo total oferecido para a sua resolução ( Questão 3 ), os estudantes apontaram, com maior incidência, a alternativa que considerava a extensão adequada, para todas as agregações consideradas ( Gráfico 4.5, Gráfico 4.6, e, no Anexo II, a Tabela II.3 ).
O percentual de alunos que responderam ser a extensão da prova adequada foi de 54,7 \%.
Já 40,5 \% dos inscritos presentes consideraram que a prova foi longa ou muito longa e menos do que 5,0 \% a avaliaram como curta ou muito curta.
Entre as Grandes Regiões, a proporção daqueles que avaliaram a prova como longa ou muito longa em relação ao tempo total destinado à sua resolução variou de 29,4 \% na região Norte até 44,5 \% na região Sudeste.
A diferença entre as regiões Norte e Sudeste é estatisticamente significativa.
Considerando -se o desempenho dos alunos, nota -se ainda que 50,1 \% consideraram a extensão da prova adequada no quarto de desempenho inferior e 61,0 \% no de melhor desempenho ( quarto superior ).
Nos quartos intermediários, esta proporção foi 55,2 \% no segundo quarto e 52,5 \% no terceiro.
No Gráfico 4.6, pode -se constatar que há uma tendência decrescente da proporção de estudantes que consideraram a prova longa ou muito longa em relação aos quartos desempenho inferior e superior : 45,7 \% e 35,5 \%, respectivamente.
Observa -se que não há diferença estatisticamente significativa entre as proporções de escolha dessa resposta nos quartos de desempenho. 4.3 COMPREENSÃO DOS ENUNCIADOS DAS QUESTÕES 4.3.1 Componente de Formação Geral Com relação aos enunciados das questões do Componente de Formação Geral ( Questão 4 ), as opiniões foram positivas, já que 79,8 \% dos alunos avaliados consideraram todos ou a maioria dos enunciados das questões “ claros e objetivos ” ( Gráfico 4.7, Gráfico 4.8, e, no Anexo II, a Tabela II.4 ).
Na análise regional, a percentagem de estudantes que avaliaram que todos ou a maioria dos enunciados das questões do Componente de Formação Geral estavam “ claros e objetivos ” variou de 76,7 \% na região Nordeste a 82,4 \% na região Sudeste, não sendo essa diferença estatisticamente significativa, assim como a diferença entre cada uma das regiões.
A análise das percepções dos estudantes sobre a clareza e objetividade dos enunciados permite afirmar que todos, ou a maioria dos enunciados de questões relativas ao Componente de Formação Geral, foram considerados “ claros e objetivos ” para a maior parte dos respondentes ( maior ou igual a 76,7 \% em todas as regiões e maior ou igual a 73,6 \% para todos os quartos de desempenho ).
Segundo o desempenho, observa -se que a proporção dos que emitiram essa opinião cresce conforme o desempenho aumenta : 73,6 \% do quarto inferior a 87,1 \% do superior.
Há diferenças estatisticamente significativas do primeiro e do segundo quartos para o quarto superior.
Nesse último quarto, a clareza e objetividade de todos ou da maioria dos enunciados das questões foi percebida por 87,1 \% dos alunos sem que seja estatisticamente significativa a diferença deste para o quarto imediatamente inferior, onde a proporção foi de 80,4 \%. 4.3.2 Componente de Conhecimento Específico Com relação aos enunciados das questões do Componente de Conhecimento Específico da prova, para 80,6 \% dos estudantes avaliados da Área de Secretariado Executivo, a clareza e a objetividade ( Questão 5 ) estavam presentes em todas ou na maioria das questões ( Gráfico 4.9, Gráfico 4.10, e no Anexo II, a Tabela II.5 ).
A maioria dos estudantes de todas as Grandes Regiões brasileiras considerou “ claros e objetivos ” todos ou a maioria dos enunciados das questões do Componente de Conhecimento Específico da prova, percentual sempre maior ou igual a 76,8 \%.
A diferença entre a maior e a menor proporção, 83,5\%(Norte) e 76,8 \% ( Nordeste ) não é estatisticamente significativa, assim como a diferença entre cada uma das regiões.
A proporção de estudantes que consideraram os enunciados das questões “ claros e objetivos ” apresenta uma tendência crescente em relação ao aumento de desempenho : mais elevada no quarto superior ( 86,0 \% ) se comparada ao quarto inferior de desempenho ( 74,2 \% ).
A diferença do primeiro para o último quarto é estatisticamente significativa. 4.4 SUFICIÊNCIA DAS INFORMAÇÕES/INSTRUÇÕES FORNECIDAS Ao avaliarem as informações/instruções fornecidas para a resolução das questões ( Questão 6 ), 88,6 \% dos respondentes da Área de Secretariado Executivo de todo o Brasil afirmaram que estas eram até excessivas ou “ suficientes ” em todas ou na maioria das questões ( Gráfico 4.11, Gráfico 4.12, e, no Anexo II, a Tabela II.6 ).
Quanto à distribuição de respondentes pelas Grandes Regiões observa -se que a proporção de estudantes que consideraram as informações/instruções fornecidas até excessivas ou “ suficientes ” em todas ou na maioria das questões foi sempre igual ou superior a 86,9 \%, chegando a 90,4 \% na região Sul.
A diferença entre o maior percentual, obtido na região Sul ( 90,4 \% ), e o menor, da região Nordeste ( 86,9 \% ) não é estatisticamente significativa, assim como a diferença entre cada uma das regiões.
Levando -se em conta o desempenho dos participantes, notam -se diferenças estatisticamente significativas entre as opiniões de estudantes do quarto inferior e superior de desempenho, como mostra o Gráfico 4.12.
O percentual de participantes que avaliaram “ as informações/instruções ” como até excessivas ou “ suficientes ” em todas ou na maioria das questões foi mais elevado no quarto superior ( 93,4 \% ), percentual superior à média nacional ( 88,6 \% ).
Já no quarto inferior, a suficiência das informações/instruções declarada como até excessiva, em todas ou na maioria das questões foi percebida por 82,9 \% dos respondentes.
Essas proporções são crescentes com os quartos de desempenho. 4.5 DIFICULDADE ENCONTRADA AO RESPONDER À PROVA Perguntados sobre as dificuldades com as quais se depararam ao responder à prova ( Questão 7 ), 11,6 \% dos estudantes apontaram o Desconhecimento do conteúdo.
Para 40,9 \%, a Forma diferente de abordagem do conteúdo foi indicada como dificuldade.
Já a Falta de motivação para fazer a prova foi a dificuldade apontada por 13,6 \% dos respondentes.
Considerando -se todo o Brasil, 29,7 \% dos respondentes afirmaram que não tiveram qualquer tipo de dificuldade para responder à prova ( Tabela II.7 no Anexo II ).
Os Gráficos 4.13 e 4.14 apresentam os percentuais de estudantes que apontaram o Desconhecimento do conteúdo como dificuldade percebida ao responder à prova.
Na análise por Grandes Regiões, o percentual de inscritos e presentes que apontaram o Desconhecimento do conteúdo como dificuldade ao responder à prova não superou 14,9 \%.
Os percentuais variaram de 8,3 \% na região Norte a 14,9 \% na região Sudeste, não sendo essa diferença estatisticamente significativa, assim como a diferença entre cada uma das regiões.
A Forma diferente de abordagem do conteúdo foi a escolha modal dos estudantes, com percentuais que variaram de 39,3 \% ( regiões Nordeste e Sudeste ) a 45,5 \% ( Centro-Oeste ).
O percentual de alunos que citaram a Falta de motivação como dificuldade variou de 8,3 \% ( região Norte ) a 17,3 \% ( região Nordeste ).
Os que declararam não ter qualquer tipo de dificuldade para responder à prova variaram de 24,1 \% na região Nordeste a 34,5 \% na região Norte.
Com relação aos quartos de desempenho, o Desconhecimento do conteúdo foi a opção escolhida por 14,4 \% dos estudantes do quarto superior e 11,5 \% do quarto inferior.
A alternativa modal para os alunos, quando agregados pelos quartos de desempenho, foi que a dificuldade encontrada foi causada pela Forma diferente de abordagem do conteúdo : 41,5 \% no quarto inferior e 34,2 \% no quarto superior assim o responderam. 4.6 CONTEÚDOS DAS QUESTÕES OBJETIVAS DA PROVA Ao analisarem os conteúdos das questões objetivas da prova ( Questão 8 ), um percentual muito pequeno dos estudantes avaliados, apenas 3,4 \%, afirmou que não estudou ainda a maioria desses conteúdos ( Gráficos 4.15, Gráfico 4.16, e a Tabelas II.8 no Anexo II ).
A maioria ( 79,2 \% ) afirmou ter estudado e aprendido muitos ou todos os conteúdos.
Na análise por Grande Região, a proporção de respondentes que escolheram a opção não estudou ainda a maioria desses conteúdos foi pequena.
Observa -se que nas regiões Nordeste ( 5,1 \% ) e Centro-Oeste ( 4,0 \% ), apesar de pequenas, as proporções foram maiores do que a média nacional ( 3,4 \% ).
Observa -se que não há diferença estatisticamente significativa entre todas as regiões.
Em todas as regiões, a maioria dos presentes afirmou ter estudado e aprendido muitos ou todos os conteúdos, com proporções variando entre 69,9 \% na região Norte e 84,7 \% na Sul.
Considerando -se separadamente as opiniões de estudantes dos quatro quartos de desempenho, observa -se que, no quarto inferior, 7,1 \% ofereceram como resposta que não estudou ainda a maioria desses conteúdos, sendo 1,1 \% os do quarto superior com a mesma resposta.
A diferença entre os alunos que optaram por esse motivo de dificuldade nos quartos extremos é estatisticamente significativa.
Tendo em conta o quarto superior, 88,5 \% dos alunos afirmaram ter estudado e aprendido muitos ou todos os conteúdos.
No outro extremo, no primeiro quarto, 67,4 \% optaram pelas mesmas categorias. 4.7 TEMPO GASTO PARA CONCLUIR A PROVA Ao responderem sobre o tempo de conclusão da prova ( Questão 9 ), mais da metade dos estudantes ( 74,6 \% ) afirmou ter gasto Entre duas e quatro horas ( Gráfico 4.17, Gráfico 4.18 e, no Anexo II, a Tabela II.9 ).
Considerando -se as cinco Grandes Regiões brasileiras, o percentual dos que utilizaram Entre duas e quatro horas para finalizar a prova apenas na região Nordeste ( 72,3 \% ) foi inferior ao percentual nacional.
Nas demais Grandes Regiões, o percentual de alunos que dispensaram Entre duas e quatro horas para concluir a prova variou de 74,8 \% ( região Centro-Oeste ) a 76,1 \% ( região Sul ), como mostra o Gráfico 4.17.
Uma vez sendo analisadas as alternativas escolhidas pelos estudantes que se situam nos diferentes quartos de desempenho, observa -se uma tendência crescente : uma maior proporção de participantes no quarto superior declarou ter gasto Entre duas e quatro horas “ para concluir a prova ” quando comparadas com os dos quartos inferiores.
As diferenças entre o primeiro ( 66,5 \% ) e o terceiro ( 79,4 \% ) e ainda entre o primeiro e o quarto de maior desempenho ( 79,5 \% ) são estatisticamente significativas.

% matched lemmas: 
\end{textsample}
